\chapter{树、概率模型与集成学习}
\label{ch:trees}

树把输入空间切成若干区域；朴素贝叶斯通过条件分布计算类别后验；集成学习组合多个模型。
本章先理解单个预测怎样产生，再解释组合的价值与限制。

\begin{keyidea}[预测结构决定了检查重点]
树关注切分与叶子，概率模型关注分布假设和概率方向，集成方法关注模型之间的依赖与组合规则。
任何一种训练方法都仍需正确划分、训练内调参和独立评估。
\end{keyidea}

\section{分类树：按规则进入叶子}
\label{sec:tree-classification}
决策树（Decision Tree）依次询问特征条件，把输入送到一个叶节点。
叶子可给出多数类别或根据训练记录估计的类别比例。
下面用八条样本说明一次二元切分。

\begin{figure}[htbp]
\centering\begin{tikzpicture}[node distance=16mm and 12mm]
\node[box,text width=45mm] (root) {$x_1\leq2.5$？\\正4 / 负4};
\node[box,text width=35mm,below left=of root] (left) {正3 / 负0\\预测正类};
\node[box,text width=35mm,below right=of root] (right) {正1 / 负4\\预测负类};
\draw[flow] (root)--node[left,lab]{是}(left);
\draw[flow] (root)--node[right,lab]{否}(right);
\end{tikzpicture}

\caption{根节点正负各4条；左叶3条，右叶5条。条件按输入特征判断。}
\par\smallskip
\readfigure{真实标签为$[1,1,1,0,0,0,0,1]$、输入$x_1=0,\ldots,7$时，
阈值2.5产生图中计数。右叶仍有一个正样本，所以叶子不一定纯，也不一定全部预测正确。}
\end{figure}

分类节点的Gini不纯度为$G=1-\sum_kp_k^2$；熵为$H=-\sum_kp_k\log_2p_k$，约定$0\log0=0$。
切分后按子节点样本数加权。此例父Gini为0.5，子节点加权Gini为
$\frac38\times0+\frac58\times[1-(1/5)^2-(4/5)^2]=0.2$，下降0.3。

\quick{训练时用标签评价候选切分；预测时沿已确定规则走树，不需要新样本的标签。}

\section{回归树：叶子输出常数}
\label{sec:tree-regression}
平方损失回归树在区域$R_m$输出训练目标均值$c_m$，
可写成$\hat f(\vct x)=\sum_mc_m\mathbf1[\vct x\in R_m]$。
每次尝试特征和阈值，使左右子区域内的残差平方和较小。

\begin{figure}[htbp]
\centering\begin{tikzpicture}
\begin{axis}[width=124mm,height=52mm,xmin=-.3,xmax=3.4,ymin=0,ymax=8,
 xlabel={$x$},ylabel={$y$},xtick={0,1,1.5,2,3},ytick={1,5,6,7},axis lines=left]
\addplot[Teal,thick] coordinates {(-.2,1)(1.5,1)};
\addplot[Teal,thick] coordinates {(1.5,6)(3.2,6)};
\addplot[Quiet,dashed] coordinates {(1.5,0)(1.5,8)};
\addplot[only marks,mark=*,Amber,mark size=3pt] coordinates {(0,1)(1,1)(2,5)(3,7)};
\node[lab] at (axis cs:.5,2) {叶均值1};
\node[lab] at (axis cs:2.5,4) {叶均值6};
\end{axis}
\end{tikzpicture}

\caption{四条记录$(0,1),(1,1),(2,5),(3,7)$，在$x=1.5$切分成两个叶子。}
\par\smallskip
\readfigure{左叶预测1，右叶预测6，平方误差和为0+1+1=2。
不切分时输出总体均值3.5，平方误差和为27。
单次切分降低训练误差，不等于已经证明泛化改善。}
\end{figure}

叶子均值是在固定区域下平方误差的最优常数；若改为绝对误差，最优常数相应为中位数。
常规轴向切分使预测呈分段常数，通常不能像线性趋势那样向训练目标范围之外外推。

\begin{pitfall}[贪心不等于全局最优树]
逐节点选择当前最佳切分是局部搜索。
它不保证给定深度下找到所有可能树中的最优结构，样本的小变化也可能改变切分。
\end{pitfall}

\section{树的复杂度与实现边界}
\label{sec:tree-control}
树越深、叶子越小，通常越能适应训练样本细节。
限制深度、叶子最少样本数、最小增益或作代价复杂度剪枝，都是控制方式。
具体参数应在训练内验证，而不是按训练误差一直生长到零。

\begin{center}
\begin{tabularx}{\textwidth}{P{25mm}P{44mm}Y}
\toprule
经典方法 & 常见选择准则 & 理解边界\\
\midrule
ID3 & 信息增益 & 经典形式面向离散分类特征。\\
C4.5 & 增益率及相关启发式 & 用于减轻取值数偏好，非简单任意最大化比值。\\
CART & 分类常用Gini，回归常用平方误差 & 二元树；可有不同损失与实现扩展。\\
\bottomrule
\end{tabularx}
\end{center}

原生类别与缺失支持取决于算法和软件版本，不能仅凭“决策树”这个名称假定都支持。
常规阈值树一般不需要数值标准化，但无序类别的整数编码仍可能改变可表达的分组。
单棵浅树通常容易读；很深的树和大量树的组合则难以完整解释。

\quick{树可以表达特征交互；不能笼统说树忽略特征之间的所有关联。}

\section{朴素贝叶斯：给定类别后作独立近似}
\label{sec:tree-naive-bayes}
朴素贝叶斯（Naive Bayes）用类先验$P(Y=c)$及各特征的类条件分布，作近似
\[
 P(\vct x\mid Y=c)\approx\prod_jP(x_j\mid Y=c),\qquad
 P(Y=c\mid\vct x)\propto P(Y=c)\prod_jP(x_j\mid Y=c).
\]
连续特征应使用相应密度，而非把单点概率写成正数。
“朴素”假设是\textbf{给定类别后的特征条件独立}，不是要求所有特征无条件独立。

\begin{center}
\begin{tabular}{lrrr}
\toprule
观测$x_1=1,x_2=0$ & 类先验 & $P(x_1=1\mid c)$ & $P(x_2=0\mid c)$\\
\midrule
类别1 & 0.6 & 0.8 & 0.5\\
类别0 & 0.4 & 0.2 & 0.25\\
\bottomrule
\end{tabular}
\end{center}
未归一化分数分别为0.24与0.02；归一化得到$P(Y=1\mid\vct x)=12/13\approx0.923$。
0--1损失下选择最大后验；错误成本不同则应比较期望成本。

\begin{pitfall}[不要在后验上再乘一次先验]
先验已经进入贝叶斯公式。特征强相关时，独立近似可能重复累计证据，
分类仍可能有效，但输出概率容易过度自信，需要独立检查。
\end{pitfall}

\section{概率模型的估计与数值实现}
\label{sec:tree-probability-estimation}
不同朴素贝叶斯版本对输入作不同假设。

\begin{center}
\begin{tabularx}{\textwidth}{P{31mm}P{45mm}Y}
\toprule
版本 & 常见输入/假设 & 需要检查\\
\midrule
Gaussian NB & 类别内数值特征的高斯密度。 & 方差与分布形态，零方差的稳定处理。\\
Multinomial NB & 词频等非负计数特征。 & 词典、平滑、非负输入；不是任意有符号数值。\\
Bernoulli NB & 二元是否出现。 & 出现与未出现都进入相应概率。\\
\bottomrule
\end{tabularx}
\end{center}

对离散计数，平滑可避免训练中未见事件使整项乘积为零。
例如有$K$种取值、总计数$N$、某取值计数$N_j$，加法平滑为$(N_j+\alpha)/(N+K\alpha)$，$\alpha>0$。
实际常用$\log P(c)+\sum_j\log P(x_j\mid c)$计算，再用稳定方式归一化。

\subsection{MLE、MAP和后验均值不能混用}
最大似然估计（Maximum Likelihood Estimation, MLE）最大化似然；
最大后验估计（Maximum a Posteriori, MAP）最大化后验。
平方损失下的贝叶斯参数点估计为后验均值。
例如硬币2次正面、1次反面，正面概率$\theta$的MLE为$2/3$。
均匀先验下，后验密度与$\theta^2(1-\theta)$成比例，即Beta$(3,2)$分布；
MAP为$2/3$，后验均值为$3/5$。均匀先验没有让均值也自动等于MLE。

\quick{参数估计的不确定性、模型的预测概率和最终决策规则，是相连但不同的对象。}

\section{Bagging与随机森林：让平均更有价值}
\label{sec:tree-bagging}
Bagging（Bootstrap Aggregating）在多份有放回重采样数据上拟合模型，再平均或投票。
随机森林（Random Forest）进一步在\textbf{每个节点的切分}中随机考虑特征子集，
常结合Bootstrap，以减少树之间相似的错误。

\begin{figure}[htbp]
\centering\begin{tikzpicture}[node distance=6mm and 10mm]
\node[box,text width=22mm] (data) {训练数据};
\node[box,text width=28mm,right=of data] (m2) {重采样2\\模型2};
\node[box,text width=28mm,above=of m2] (m1) {重采样1\\模型1};
\node[box,text width=28mm,below=of m2] (m3) {重采样3\\模型3};
\node[box,text width=29mm,right=16mm of m2] (agg) {平均或投票\\组合预测};
\foreach \m in {m1,m2,m3}{\draw[flow] (data.east)--(\m.west);\draw[flow] (\m.east)--(agg.west);}
\node[lab,below=6mm of m3] {各模型可分别训练，预测误差仍可能相关};
\end{tikzpicture}

\caption{并列训练与组合示意；仅画三个模型以说明流程。}
\par\smallskip
\readfigure{回归可平均预测。分类可硬投票或平均类别概率后决策，具体要看算法实现。
三模型正类概率0.49、0.49、0.99，硬投票预测0，但概率均值约0.657会预测1。}
\end{figure}

若$B$个模型预测方差都为$\sigma^2$、两两相关系数为$\rho$，其均值方差为
\[
 \sigma^2\left[\rho+\frac{1-\rho}{B}\right].
\]
$\sigma^2=1,\rho=0.2$时，10模型得0.28，100模型得0.208。
这个简化分析说明相关误差不能仅靠无限增加树数消除，不是对所有森林的精确测量公式。

袋外（Out-of-Bag）记录可用于相应评估，但分组、时间约束与数据处理也需保持正确。
随机森林能降低单树的高方差风险，仍可能过拟合或遭遇分布变化。

\section{袋外估计：只汇总没见过该样本的树}
\label{sec:tree-oob}
对$n$条记录有放回抽取$n$次构成一棵树的Bootstrap训练集。
某条记录可能被抽到多次，也可能一次都没抽到；后一种记录对这棵树称为袋外
（Out-of-Bag, OOB）。它不是对整个森林都未参与训练的独立测试记录。

\begin{figure}[htbp]
\centering\begin{tikzpicture}[x=13mm,y=8mm]
\foreach \x in {1,...,5}{\node[lab] at (\x,4.7) {样本\x};}
\foreach \y/\t in {4/1,3/2,2/3,1/4}{\node[lab,anchor=east] at (.3,\y) {树\t};}
\foreach \x/\y/\v in {1/4/2,2/4/1,4/4/2,2/3/1,3/3/2,5/3/2,1/2/1,2/2/2,3/2/1,5/2/1,1/1/1,3/1/1,4/1/2,5/1/1}{
 \node[draw=Teal,fill=Pale,minimum width=12mm,minimum height=7mm,font=\small] at (\x,\y) {\v};}
\foreach \x/\y in {3/4,5/4,1/3,4/3,4/2,2/1}{
 \node[draw=Amber,fill=Amber!15,minimum width=12mm,minimum height=7mm,font=\small] at (\x,\y) {0};}
\node[box,text width=43mm] at (8,2.5) {样本4的袋外预测\\仅来自树2、树3\\$(0.8+0.6)/2=0.7$};
\draw[Amber,thick] (4.42,1.55)--(4.5,1.55)--(4.5,3.45)--(4.42,3.45);
\draw[flow] (4.5,2.5)--(6.1,2.5);
\end{tikzpicture}

\caption{5条原记录、4份Bootstrap样本。格中为被抽次数；橙色0表示袋外。每行之和为5。}
\par\smallskip
\readfigure{对样本4，只有树2与树3没有用它训练。
若两树对该样本给正类概率0.8、0.6，袋外平均为0.7；树1、树4的预测不参与这一次袋外汇总。}
\end{figure}

每次未抽中某条记录的概率为$1-1/n$，独立抽$n$次后
\[
 P(\text{该记录袋外})=\left(1-\frac1n\right)^n\longrightarrow e^{-1}\approx0.368.
\]
$n=5$时是$(4/5)^5=0.32768$，不是精确36.8\%。
期望袋外记录数为$n(1-1/n)^n$；同一Bootstrap样本中不同记录的袋外指示并不独立。

\begin{enumerate}
\item 为每条记录找出没有抽到它的树，汇总这些树的预测。
\item 将袋外预测与该记录真实标签比较，再按指定指标汇总记录。
\item 检查袋外树数量；树太少时，部分记录可能没有袋外预测或估计很不稳定。
\end{enumerate}
不是先给每棵树算一个误差，再默认这些误差的简单均值就是同一个袋外指标。
分类平均概率还是硬投票，也需与实现约定一致。

\begin{pitfall}[袋外不自动评估整套预处理]
若特征筛选已利用全部训练标签，袋外记录的信息仍可能提前进入树。
袋外评分不能自动修复这种泄漏。评估完整流水线时，可采用折内拟合的交叉验证。
反复按袋外分数选参也会产生选择影响；分组和时间数据不能只靠逐行重采样模拟部署情形。
\end{pitfall}


\section{特征重要性：量的是哪种贡献？}
\label{sec:tree-importance}
树的不纯度重要性（Mean Decrease in Impurity, MDI）累计某特征参与切分时带来的加权不纯度下降。
它反映训练中的使用方式，可能偏向可尝试切分较多的高基数特征。

\begin{center}
\begin{tabularx}{\textwidth}{P{35mm}Y}
\toprule
方法 & 解释与限制\\
\midrule
不纯度重要性 & 来自训练切分；高基数、相关特征和过拟合可影响排序。\\
置换重要性 & 固定已拟合模型，打乱特征后比较评估分数变化；依赖数据、指标与模型。\\
分组重要性 & 可把同一原始变量的多个编码列作为一组处理。\\
\bottomrule
\end{tabularx}
\end{center}

若一个类别变量拆成三个独热列，MDI分别为0.10、0.04、0.06，
三者求和0.20可作为该模型对原变量的汇总，但不自动消除基数偏好。
若做分组置换，应对整组使用同一行排列，以保持合法的独热组合；不能随意独立打乱各列。

置换重要性可在开发数据的留出部分计算。若据此反复选择特征或模型，该部分已经参与选择，
不能再把其结果当作完全未使用的最终测试。

\begin{pitfall}[相关特征可能掩盖彼此的重要性]
一个特征打乱后，另一个仍提供相似信息，分数下降可能很小。
这不等于两者都无信息；高重要性也不证明因果作用。
\end{pitfall}

\section{梯度提升：逐步拟合修正方向}
\label{sec:tree-gradient-boosting}
梯度提升（Gradient Boosting）构建加法模型$F_m=F_{m-1}+\nu h_m$。
新模型拟合当前损失对预测值的负梯度方向；$\nu>0$是收缩步长。
使用$\ell(y,F)=\frac12(y-F)^2$时，这个方向恰好是残差$y_i-F_{m-1}(x_i)$；
若省略$1/2$，负梯度为残差的2倍，需相应理解步长的尺度。

\begin{figure}[htbp]
\centering\begin{tikzpicture}[node distance=8mm]
\node[box,text width=25mm] (pred) {当前预测\\$F_{m-1}$};
\node[box,text width=29mm,right=of pred] (grad) {计算修正目标\\负梯度};
\node[box,text width=25mm,right=of grad] (tree) {拟合新模型\\$h_m$};
\node[box,text width=29mm,right=of tree] (new) {加到已有预测\\$F_m=F_{m-1}+\nu h_m$};
\draw[flow] (pred)--(grad);\draw[flow] (grad)--(tree);\draw[flow] (tree)--(new);
\draw[flow] (new.south)--++(0,-10mm)-|(pred.south);
\node[lab,below=15mm of grad] {下一轮依赖当前组合的预测};
\end{tikzpicture}

\caption{每轮依赖当前预测。一般损失的修正目标不一定是普通数值残差。}
\par\smallskip
\readfigure{目标$[0,2,4]$，初始常数预测为2。
残差$[-2,0,2]$若由一个树桩拟合成$[-2,1,1]$，取$\nu=0.5$，
新预测为$[1,2.5,2.5]$，训练MSE从$8/3$下降为$7/6$。}
\end{figure}

Boosting不都通过“给错样本加权”实现；梯度提升也不是只拟合前一棵树的错误，
而是针对\textbf{当前全部组合}的损失更新。树深、学习率、轮数、采样与正则化共同影响结果。
XGBoost、LightGBM等是相关的树提升实现，细节和默认行为应查具体版本文档。

\quick{训练损失下降是一个过程信号；轮数与停止仍应根据合适的验证表现选择。}

\section{AdaBoost：一个完整的二分类权重步骤}
\label{sec:tree-adaboost}
经典离散二分类AdaBoost使用$y_i,h_m(x_i)\in\{-1,1\}$及归一化样本权重$D_m(i)$。
\[
 e_m=\sum_iD_m(i)\mathbf1[h_m(x_i)\ne y_i],\qquad
 \alpha_m=\frac12\log\frac{1-e_m}{e_m},
\]
\[
 D_{m+1}(i)=\frac{D_m(i)\exp[-\alpha_my_ih_m(x_i)]}{Z_m},\qquad
 H(x)=\operatorname{sign}\!\left(\sum_m\alpha_mh_m(x)\right).
\]
$Z_m$使权重和为1；边界、停止与平票需约定。

\begin{figure}[htbp]
\centering\begin{tikzpicture}
\begin{axis}[width=125mm,height=48mm,ymin=0,ymax=.58,xmin=.5,xmax=4.5,
 xlabel={训练样本},ylabel={归一化权重},xtick={1,2,3,4},ytick={0,.1666667,.25,.5},
 yticklabels={0,$1/6$,$1/4$,$1/2$},ybar,bar width=7mm,
 legend style={font=\small,at={(.5,1.03)},anchor=south,legend columns=2}]
\addplot[fill=Teal!60,draw=Teal] coordinates {(1,.25)(2,.25)(3,.25)(4,.25)};\addlegendentry{更新前}
\addplot[fill=Amber!60,draw=Amber] coordinates {(1,.1666667)(2,.1666667)(3,.1666667)(4,.5)};\addlegendentry{更新后}
\end{axis}
\end{tikzpicture}

\caption{四个样本初始均为1/4，仅第4个被当前弱模型分错。}
\par\smallskip
\readfigure{$e=1/4$、$\alpha=\frac12\log3\approx0.549$。
更新归一化后，错样本权重为1/2，另外三个各为1/6；
这是相对权重变化，不是重新采集了更多错样本。}
\end{figure}

若$e=0$，应采用相应提前停止/有限数值处理；不能直接除零。
若$e\geq1/2$，没有上述正权重弱优势，需按算法策略停止或处理，而非机械继续。
噪声标签和离群记录可能被反复强调，轮数也需验证。

\section{投票、Boosting与Stacking怎样选择？}
\label{sec:tree-compare}
集成不是保证每个模型都正确，而是利用预测之间的差异及组合结构。
如果所有模型犯同样的错误，增加数量帮助有限。

\begin{center}
\begin{tabularx}{\textwidth}{P{28mm}P{49mm}Y}
\toprule
方法 & 模型如何形成 & 常见检查\\
\midrule
Bagging/森林 & 重采样与随机化后分别拟合。 & 相关性、叶子限制、袋外/独立验证。\\
Boosting & 根据已有组合逐步添加修正模型。 & 学习率与轮数、损失、噪声敏感性。\\
Voting/平均 & 组合多个已定义的候选模型。 & 权重、概率与硬标签组合的差别。\\
Stacking & 用基础模型的预测训练一个组合模型。 & 组合模型需使用折外预测，防止训练预测泄漏。\\
\bottomrule
\end{tabularx}
\end{center}

Stacking中，基础模型在某样本上的输入给第二层的预测，通常应由未用该样本训练的模型产生。
用第一层对自己训练样本的拟合预测直接训练第二层，会让第二层看到过于乐观的输入。
最终仍要评估包含预处理、第一层和第二层的完整流程。

\subsection{自测与回看}
森林的树是否必然相互独立？\textbf{不是。}
梯度提升是否总是修改样本权重？\textbf{不是。}
贝叶斯后验上是否还要再乘先验？\textbf{不需要。}
特征重要性是否等于因果作用？\textbf{不等于。}

\begingroup\small
\noindent\textbf{参考：}ISLR第二版第4、8章；scikit-learn的Decision Trees、Naive Bayes、Ensembles和Permutation importance文档。
数字、流程图和权重图均独立构造。
\par\endgroup
